\clearpage
\section{After Five Turns}

\begin{lessonbox}
The example stops at a natural transition point. The Partisans have organized
into brigades, German movement has expanded, the Chetniks are politically
fragmented, and Tito is temporarily off map. Game-Turn 6 will introduce Allied
Progress and a new Weather determination.
\end{lessonbox}

\begin{multicols}{2}
\small
\subsection{End Position}

\rowcolors{2}{TitoBlue!42}{white}
\begin{tabularx}{\linewidth}{@{}>{\bfseries}l X@{}}
  \rowcolor{TitoInk!15}
  Item & State after Game-Turn 5 \\
  Yugoslav VP & 24 \\
  Partisans & 30 groups, 2 brigades \\
  Tito & Withdrawn; returns GT9 \\
  Chetniks & 2 pro-Yugoslav, 4 pro-Axis, 1 neutral \\
  AGO reserve & 5 remain from the secret limit of 8 \\
  Weather & Drought ends with GT5 \\
  Brigade trigger & In effect since GT3 \\
  Bulgarian 25/27 & Deployed in Serbia \\
  German 342 & Transferred on GT5 \\
\end{tabularx}
\rowcolors{2}{}{}

\subsection{What Happens on Game-Turn 6?}

\begin{enumerate}[leftmargin=1.4em,itemsep=2pt]
  \item Begin Allied Progress in Box 1 and make the first progress roll
    \rulesref{10.11}.
  \item Make the GT6 Weather roll, establishing conditions through GT9
    \rulesref{11.22B}.
  \item Tito remains off map; skip identification and location.
  \item Deploy German 7SS \rulesref{13.91}.
  \item Continue collaboration and decide whether another AGO is worth one of
    the five remaining operations.
\end{enumerate}

\begin{commentarybox}[title={Do not read 24 VP as a forecast}]
The tutorial makes conservative Objective choices and deliberately exposes
Tito. Its running total exists to teach the Victory Point Stage, not to model
expert Yugoslav play or predict the final result.
\end{commentarybox}

\columnbreak

\subsection{Five Opening Lessons}

\begin{enumerate}[leftmargin=1.4em,itemsep=3pt]
  \item \textbf{Occupation is temporary.} Guerrillas can collect VP and
    recruits without defeating every Axis unit, but an Axis response can force
    them back into the Hide-aways.
  \item \textbf{The political map matters.} Chetnik allegiance changes can
    transfer whole stacks between players or remove them from both players'
    control.
  \item \textbf{Reaction is the defense against AGOs.} Moving Tito or an
    organized unit before deployment may matter more than preserving isolated
    groups.
  \item \textbf{Garrison arithmetic drives growth.} A large Zone shortfall can
    create more Partisans than routine Objective recruitment.
  \item \textbf{Triggers reshape the war.} The first brigade widens German
    movement, releases Bulgarian reinforcements, and starts the hunt for Tito.
\end{enumerate}

\subsection{Rules Used Most Often}

\begin{tabularx}{\linewidth}{@{}>{\bfseries}l X@{}}
  4.2 & Sequence of Play \\
  6.1--6.4 & Movement and stacking \\
  7.2 & Chetnik collaboration \\
  7.4--7.6 & Recruitment, uprisings, organization \\
  8.1--8.3 & Combat and terrain \\
  8.4 & Anti-Guerrilla Operations \\
  9.3--9.4 & Tito, recruitment, and the hunt \\
  11.0 & Weather \\
  13.9 & Reinforcement and transfer schedule \\
  14.2--14.3 & Victory Points \\
\end{tabularx}
\end{multicols}

\vfill
{\color{TitoInk}\hrule height 0.7pt}
\vspace{5pt}
\begin{center}
  {\sffamily\bfseries\large Continue the game from Game-Turn 6—or reset the
  pieces and try a different opening.\par}
  \vspace{5pt}
  \counterpic[0.58in]{tito-identified}\hspace{0.14in}
  \counterpic[0.58in]{partisan-brigade}\hspace{0.14in}
  \counterpic[0.58in]{german-parachute}\hspace{0.14in}
  \counterpic[0.58in]{anti-guerrilla-operations}
\end{center}
